\documentclass[../../../../main.tex]{subfiles}
\begin{document}

$n$を2より大きい正の整数とする．
曲線
\begin{align*}
  y=x^n \hspace{1zw} (i)
\end{align*}
上で，$x$座標が0，1，2である点をそれぞれ$O$，$A$，$B$とし，
$O$，$A$，$B$をとおり$y$軸に平行な軸をもつ放物線
\begin{align*}
  y=f(x) \hspace{1zw} (ii)
\end{align*}
をえがく．
曲線(i)および放物線(ii)の，
$O$，$A$の間にある部分の囲む面積を$S_1$，
$A$，$B$の間にある部分の囲む面積を$S_2$とするとき，
$S_1=S_2$となるためには，$n$はどのような数でなければならないか．

\end{document}
